% Part: first-order-logic
% Chapter: introduction
% Section: sentences

\documentclass[../../../include/open-logic-section]{subfiles}

\begin{document}

\olfileid{fol}{int}{snt}

\olsection{\usetoken{P}{sentence}}

چنگا، ہن !!{structure}s تے !!{formula}s
لئی پورا اترن (یعنی ”وچ سچا ہون“)
دی تعریف دا اک خاکہ ساڈے کول اے۔
پر اوہنوں اک وادھو شے---!!a{variable}
قدر-بندی---دی لوڑ اے، جدوں کہ اسیں
!!{sentence}s دی تعریف چاہندے ساں۔
قدر-بندیاں توں چھٹکارا کیویں
پائیے، تے !!{sentence}s اصل وچ
کیہ نیں؟

رسمی منطق دے اپنے پہلے تعارف وچ
تہانوں شاید !!{variable}s تے
مسوّراں دے رشتے بارے بحث یاد
ہووے۔ مسوّر دے فوراً بعد
ہمیشہ !!a{variable} آؤندا اے؛
فیر !!{sentence} دے اوس حصے
وچ جتھے مسوّر لاگو ہوندا اے
(اوہدے ”دائرۂ اثر“ وچ)،
اوہ !!{variable} اوس مسوّر
نال ”پابند“ سمجھیا جاندا اے۔
!!{formula}s وچ ایہ ضروری
نہیں سی کہ ہر !!{variable}
لئی کوئی متعلقہ مسوّر وی ہووے؛
جیہڑے !!{variable}s دے
متعلقہ مسوّر نہ ہون، اوہ
”آزاد“ یا ”غیر پابند“ ہوندے
نیں۔ اسیں !!{sentence}s
اوہناں سارے !!{formula}s
نوں آکھاں گے جنہاں وچ
کوئی آزاد !!{variable}s
نہ ہون۔

اک واری فیر، !!a{formula}~$!A$
وچ !!a{variable} دی کسے
آمد بارے ایہ سمجھنا اوکھا
نہیں کہ اوہ کدوں پابند اے،
کیہڑا مسوّر اوہنوں پابند
کردا اے، تے کدوں آزاد اے۔
ممکن اے تسی ایس دی جانچ
دا کوئی طریقہ وی سکھیا
ہووے، شاید قوسین گن کے۔
پر سانوں ٹھیک تعریف
چاہیدی اے---تے چونکہ
اسیں !!{formula}s نوں
استقرائی طور تے تعریف
کیتا اے، ایس لئی
!!a{variable}~$x$
دیاں آزاد تے پابند
آمدواں دی تعریف
!!a{formula}~$!A$ وچ وی
استقرائی طور تے دے
سکدے آں۔ ساڈی سادہ
زبان لئی ایہ کجھ
ایس طرح ہو سکدی اے:
\begin{enumerate}
  \item جے $!A$ ایٹمی اے،
  تاں اوہدے وچ $x$ دیاں
  ساریاں آمدواں آزاد نیں
  (یعنی $x$ دی آمد
  ~$\Atom{\Obj P}{x}$ وچ آزاد اے)۔
  \item جے $!A$ دا روپ
  $\lnot !B$ اے، تاں
  $x$ دی کوئی آمد
  ~$\lnot !B$ وچ اُدوں
  تے صرف اُدوں آزاد
  اے جدوں اوہدے مطابق
  آمد~$x$،
  ~$!B$ وچ آزاد اے
  (یعنی~$
  !B$ وچ متغیراں
  دیاں آزاد آمدواں
  عین~$\lnot !B$
  وچ انہاں دے مطابق
  آمدواں نیں)۔
  \item جے $!A$ دا روپ
  $(!B \land !C)$ اے،
  تاں~$x$ دی کوئی آمد
  $(!B \land !C)$ وچ
  اُدوں تے صرف اُدوں
  آزاد اے جدوں اوہدے
  مطابق آمد~$x$،
  ~$!B$ یا~$!C$
  وچ آزاد اے۔
  \item جے $!A$ دا روپ
  $\lexists[x][!B]$ اے،
  تاں $x$ دی کوئی
  وی آمد~$!A$ وچ
  آزاد نہیں؛ جے روپ
  $\lexists[y][!B]$
  اے، جتھے $y$،
  ~$x$ توں وکھرا
  !!{variable} اے،
  تاں $x$ دی کوئی
  آمد $\lexists[y][!B]$
  وچ اُدوں تے صرف
  اُدوں آزاد اے جدوں
  اوہدے مطابق آمد~$x$،
  ~$!B$ وچ آزاد اے۔
\end{enumerate}

جدوں متغیراں دیاں
آزاد تے پابند آمدواں
دی ٹھیک تعریف ساڈے
کول ہووے، تاں اسیں
سادہ طور تے کہہ سکدے
آں: !!a{sentence} اوہ
ہر !!{formula} اے
جیہدے وچ !!{variable}s
دیاں آزاد آمدواں
نہ ہون۔

\end{document}
